% =====================================================
% 题型：平面向量三角形交点解答题 (q_vector_triangle_intersection.tex)
% =====================================================
% 难度：★★★ (难题)
% =====================================================
\newcommand{\qVectorCevianIntersectionStem}{%
  如图，在 \(\triangle ABC\) 中，\(AB = 2\)，\(AC = 3\)，\(\angle BAC = \dfrac{\pi}{3}\)，\(M\) 为 \(BC\) 的中点，\(2\vec{AN} = \vec{NC}\)，\(AM, BN\) 相交于点 \(P\)。
  \begin{enumerate}[label=(\arabic*) , leftmargin=2em, itemsep=0.1em]
    \item 设 \(\vec{AB} = \vec{a}\)，\(\vec{AC} = \vec{b}\)，用 \(\vec{a}\)，\(\vec{b}\) 表示 \(\vec{BP}\)；
    \item 求 \(\sin\angle BPM\)。
  \end{enumerate}%
}

\newcommand{\qVectorCevianIntersectionFigure}[1][1.5]{%
  \begin{tikzpicture}[scale=#1, >=stealth]
    % Triangle vertices
    \coordinate (A) at (0,0);
    \coordinate (C) at (2.4,0);
    \coordinate (B) at (0.8,1.386); % AB=2, AC=3, angle=60

    % Midpoint M of BC
    \coordinate (M) at ($(B)!0.5!(C)$);

    % Point N on AC such that 2AN = NC => AN = 1/3 AC
    \coordinate (N) at ($(A)!0.333!(C)$);

    % Intersection P of AM and BN
    % AP = 1/4 a + 1/4 b
    \coordinate (P) at ($(A)!0.5!(M)$);

    \draw[thick] (A) -- (B) -- (C) -- cycle;
    \draw[thick] (A) -- (M);
    \draw[thick] (B) -- (N);

    \node[below left] at (A) {\small $A$};
    \node[above] at (B) {\small $B$};
    \node[below right] at (C) {\small $C$};
    \node[right] at (M) {\small $M$};
    \node[below] at (N) {\small $N$};
    \node[above left] at (P) {\small $P$};
  \end{tikzpicture}%
}

\newcommand{\qVectorCevianIntersectionAnswer}{%
  (1) \(\vec{BP} = -\dfrac{3}{4}\vec{a} + \dfrac{1}{4}\vec{b}\)；\\
  (2) \(\sin\angle BPM = \dfrac{4\sqrt{19}}{19}\)。%
}

\newcommand{\qVectorCevianIntersectionSolution}{%
  (1) 因为 \(M\) 是 \(BC\) 的中点，且 \(\vec{AB} = \vec{a}\)，\(\vec{AC} = \vec{b}\)，
  所以 \(\vec{AM} = \dfrac{1}{2}(\vec{AB} + \vec{AC}) = \dfrac{1}{2}(\vec{a} + \vec{b})\)。
  设 \(\vec{AP} = \lambda \vec{AM} = \dfrac{\lambda}{2}(\vec{a} + \vec{b})\)。
  因为 \(2\vec{AN} = \vec{NC} \implies \vec{AN} = \dfrac{1}{3}\vec{AC} = \dfrac{1}{3}\vec{b}\)。
  由于点 \(P\) 在线段 \(BN\) 上，可设 \(\vec{AP} = \mu \vec{AB} + (1-\mu)\vec{AN} = \mu \vec{a} + \dfrac{1-\mu}{3}\vec{b}\)。
  根据平面向量基本定理，基底 \(\vec{a}, \vec{b}\) 的系数唯一对应，得：
  \(\begin{cases}
    \dfrac{\lambda}{2} = \mu \\
    \dfrac{\lambda}{2} = \dfrac{1-\mu}{3}
  \end{cases} \implies \mu = \dfrac{1-\mu}{3} \implies 3\mu = 1-\mu \implies \mu = \dfrac{1}{4}\)。
  所以 \(\lambda = 2\mu = \dfrac{1}{2}\)。
  故 \(\vec{AP} = \dfrac{1}{4}\vec{a} + \dfrac{1}{4}\vec{b}\)。
  则 \(\vec{BP} = \vec{AP} - \vec{AB} = \left(\dfrac{1}{4}\vec{a} + \dfrac{1}{4}\vec{b}\right) - \vec{a} = -\dfrac{3}{4}\vec{a} + \dfrac{1}{4}\vec{b}\)。\\

  (2) 由 (1) 知，\(\vec{BP} = -\dfrac{3}{4}\vec{a} + \dfrac{1}{4}\vec{b} \implies \vec{PB} = \dfrac{3}{4}\vec{a} - \dfrac{1}{4}\vec{b}\)。
  又因为 \(\vec{PM} = \vec{AM} - \vec{AP} = \dfrac{1}{2}(\vec{a} + \vec{b}) - \left(\dfrac{1}{4}\vec{a} + \dfrac{1}{4}\vec{b}\right) = \dfrac{1}{4}\vec{a} + \dfrac{1}{4}\vec{b}\)。
  已知 \(|\vec{a}| = 2\)，\(|\vec{b}| = 3\)，且 \(\angle BAC = \dfrac{\pi}{3}\)，
  所以 \(\vec{a} \cdot \vec{b} = |\vec{a}||\vec{b}|\cos\dfrac{\pi}{3} = 2 \times 3 \times \dfrac{1}{2} = 3\)。
  计算 \(\vec{PM}\) 的模长：
  \(|\vec{PM}|^2 = \dfrac{1}{16}|\vec{a} + \vec{b}|^2 = \dfrac{1}{16}(|\vec{a}|^2 + 2\vec{a}\cdot\vec{b} + |\vec{b}|^2) = \dfrac{1}{16}(4 + 6 + 9) = \dfrac{19}{16} \implies |\vec{PM}| = \dfrac{\sqrt{19}}{4}\)。
  计算 \(\vec{PB}\) 的模长：
  \(|\vec{PB}|^2 = \dfrac{1}{16}|3\vec{a} - \vec{b}|^2 = \dfrac{1}{16}(9|\vec{a}|^2 - 6\vec{a}\cdot\vec{b} + |\vec{b}|^2) = \dfrac{1}{16}(36 - 18 + 9) = \dfrac{27}{16} \implies |\vec{PB}| = \dfrac{3\sqrt{3}}{4}\)。
  计算 \(\vec{PM} \cdot \vec{PB}\) 的数量积：
  \(\vec{PM} \cdot \vec{PB} = \left(\dfrac{1}{4}\vec{a} + \dfrac{1}{4}\vec{b}\right) \cdot \left(\dfrac{3}{4}\vec{a} - \dfrac{1}{4}\vec{b}\right) = \dfrac{1}{16}(3|\vec{a}|^2 + 2\vec{a}\cdot\vec{b} - |\vec{b}|^2) = \dfrac{1}{16}(12 + 6 - 9) = \dfrac{9}{16}\)。
  设 \(\vec{PM}\) 与 \(\vec{PB}\) 的夹角为 \(\theta\)，则 \(\theta = \angle BPM\)：
  \(\cos\angle BPM = \dfrac{\vec{PM} \cdot \vec{PB}}{|\vec{PM}||\vec{PB}|} = \dfrac{9/16}{\dfrac{\sqrt{19}}{4} \times \dfrac{3\sqrt{3}}{4}} = \dfrac{3}{\sqrt{57}} = \dfrac{\sqrt{57}}{19}\)。
  由于 \(\angle BPM \in (0, \pi)\)，所以：
  \(\sin\angle BPM = \sqrt{1 - \cos^2\angle BPM} = \sqrt{1 - \dfrac{57}{361}} = \sqrt{\dfrac{304}{361}} = \dfrac{4\sqrt{19}}{19}\)。%
}

\newcommand{\qVectorCevianIntersectionDiagnosis}{%
  若第 (1) 问已经得到 \(\vec{AP}=\frac14\vec a+\frac14\vec b\)，但第 (2) 问又改用了不一致的交点参数，后续 \(\vec{PM}\)、数量积和夹角会整体偏离。另一个高频断点是混淆 \(\vec{BP}\) 与 \(\vec{PB}\) 的方向。%
}

\newcommand{\qVectorCevianIntersectionTeachingNotes}{%
  讲评时先用两条直线参数法稳定求交点，再强调 \(\angle BPM\) 的两边必须从顶点 \(P\) 出发，应使用 \(\vec{PB}\) 与 \(\vec{PM}\) 计算。%
}
